\section{有效位、气泡与流水寄存器的所有权}

流水寄存器不能只保存“看起来像一条指令”的数据，还要有有效位说明这些字段是否代表一项真实工作。复位、停顿、冲刷和异常都会改变有效位：复位后各级无有效指令；前端停顿时 PC 与 IF/ID 保持，但允许下游继续排空；插入气泡时下一级得到 code{valid=0}，所有写使能随之失效；冲刷则把错误路径上的年轻项标成无效。气泡不是一条编码为 NOP 的普通指令，因为普通 NOP 仍可能携带取指异常、性能计数或调试语义。

每一级对流水寄存器只有三种互斥动作：保持旧值、接收上一级的新值、写入无效项。控制器必须先按优先级处理复位和异常重定向，再处理资源等待，最后处理正常推进。若同一拍既发生 Cache miss 又发现更老指令异常，机器应保留能完成精确异常所需的老状态，同时撤销年轻请求；简单地让“stall 优先”可能阻止异常入口，简单地让“flush 清空一切”又可能丢失尚未归属的总线响应。

\begin{longtable}{L{0.18\linewidth}L{0.32\linewidth}L{0.40\linewidth}}
\toprule
事件 & 流水寄存器动作 & 不能丢失的状态 \\
\midrule
正常推进 & 下一级接收数据、控制与 valid & PC、目标寄存器、异常元数据 \\
上游停顿 & PC 与相关前级保持，下游可排空 & 被阻塞指令及其控制位 \\
插入气泡 & 下一级 valid 清零 & 老指令仍按原路径推进 \\
分支冲刷 & 错误路径各级 valid 清零，PC 重定向 & 分支自身结果与正确目标 \\
可变时延等待 & 请求级保持事务身份，响应级按 ID 接收 & 地址、byte enable、请求 ID、错误状态 \\
精确异常 & 阻止年轻副作用并保存异常 PC/原因 & 老指令提交状态和异常指令身份 \\
\bottomrule
\end{longtable}

\section{多周期功能单元与可变时延存储器}

固定五级图默认 EX 和 MEM 各在一个周期内给出结果，真实乘除法器、Cache miss、TLB miss 和设备访问却具有可变时延。流水线因此需要明确请求、接受、响应和取消边界。执行级把操作数和操作码交给乘法器，只有在请求握手成功后才能认为功能单元取得所有权；乘法器忙时，相关指令和其上游依赖项必须停顿，独立功能单元是否继续运行则取决于核心是否支持多发射或记分牌。

数据访问也不能用固定倒计时替代响应。load 请求被 Cache 接受后，流水状态要保存目标寄存器、访问宽度、符号扩展方式和异常归属。命中可以很快返回；缺失则由 MSHR 继续追踪，响应可能在若干十到数百周期后到达。若核心只允许一笔阻塞式访问，可以冻结前端直到返回；若允许多笔在途，就需要 tag、返回队列和提交顺序把响应重新归属到正确指令。

store 的边界更严格。地址和数据算出不代表可以立即改写外部状态。普通内存 store 可以先进入 store buffer，待指令达到提交点后再向 Cache 取得所有权；MMIO store 往往具有不可撤销副作用，必须在所有更老异常已经排除且协议允许时才发出。第 9 章讨论的精确异常，落实到这里就是“年轻指令不能越过不可恢复的设备副作用边界”。

\section{分支预测器保存什么状态}

没有预测时，前端只能等待分支方向与目标确定，深流水线会频繁空闲。动态预测把历史状态放在取指路径上：方向预测表（Branch History Table, BHT）用饱和计数器估计是否跳转；分支目标缓冲（Branch Target Buffer, BTB）缓存已见分支的目标；返回地址栈（Return Address Stack, RAS）预测函数返回目标。三者分别回答“跳不跳”“跳到哪”“返回到哪”，不能互相替代。

预测只选择下一条要取的指令，不改变架构状态。分支在执行级解析后比较实际方向/目标与预测结果：正确则保留年轻指令，错误则清除错误路径 valid、恢复正确 PC，并更新预测器。预测器更新通常可以在解析时发生，但其状态必须按实现定义处理异常、上下文切换和安全隔离；预测器内容不是程序可见寄存器，却会影响性能和微架构侧信道。

二位饱和计数器有强不跳转、弱不跳转、弱跳转、强跳转四态。实际跳转使状态向强跳转移动一步，未跳转使其向强不跳转移动一步。它需要连续两次相反结果才从一个强状态翻到相反预测，可避免循环末次退出立即破坏下一轮的整体预测。BTB 命中但 BHT 预测不跳转时仍取顺序 PC；BHT 预测跳转但 BTB 没有目标时，前端可能只能等待译码形成目标。

\begin{longtable}{L{0.18\linewidth}L{0.28\linewidth}L{0.34\linewidth}L{0.12\linewidth}}
\toprule
结构 & 索引/输入 & 保存状态 & 输出 \\
\midrule
BHT & 取指 PC 与可选历史 & 饱和计数器或模式状态 & 跳转方向预测 \\
BTB & 取指 PC & tag、分支类型、目标地址 & 预测目标 \\
RAS & call/return 事件 & 有限深度的返回地址栈 & 预测返回 PC \\
恢复检查点 & 预测发生时的重命名/前端状态 & 可恢复指针或快照 & 误预测后的正确状态 \\
\bottomrule
\end{longtable}

\section{逐拍核对一次相关、停顿与重定向}

考虑 code{LOAD R3,[R1]}、code{ADD R5,R3,R4}、code{BEQ R5,R0,target} 三条连续指令。假设 load 命中但数据到 MEM 末端才有效，分支在 EX 末端解析。第 1 至第 3 拍依次形成正常流水；第 4 拍 ADD 原本要进入 EX，却拿不到 R3，因此 IF/ID 保持，ID/EX 写入气泡，LOAD 继续进入 MEM。第 5 拍数据可由 MEM/WB 旁路给 ADD，流水恢复。分支随后读取经旁路得到的 R5；若预测错误，则在解析沿之后清除 IF 与 ID 中的年轻项，并把 PC 改为 target。

这条 trace 要同时检查四类状态：数据值是否来自最新生产者，valid 是否随停顿和冲刷正确变化，分支之前的老指令是否仍能提交，以及错误路径是否从未写寄存器、Cache 或 MMIO。只核对最终 R5 的数值，无法发现错误路径曾短暂写设备寄存器的严重缺陷。

\begin{keyidea}
流水线控制的核心是所有权：每项工作在哪一级、是否有效、依赖谁、何时可以产生不可撤销副作用。旁路解决“值已产生但尚未写回”，停顿解决“值尚未产生”，冲刷解决“这项工作本来就不该执行”。
\end{keyidea}

\section*{章末练习}

\begin{longtable}{L{0.18\linewidth}L{0.42\linewidth}L{0.30\linewidth}}
\toprule
任务 & 要完成的产物 & 检查要点 \\
\midrule
有效位 & 比较气泡与编码 NOP 的状态含义 & 气泡不携带真实工作或副作用 \\
优先级 & 同拍发生老指令异常和年轻 load 等待时给出控制顺序 & 精确异常优先且响应仍可归属 \\
可变时延 & 列出 load miss 等待期间必须保存的字段 & 地址、目标、宽度、ID 和错误同行 \\
预测结构 & 分别说明 BHT、BTB 与 RAS 的输入和输出 & 方向、普通目标、返回目标分开 \\
计数器 trace & 从弱跳转开始，对 T、T、N、N、N 更新二位计数器 & 每次只向实际结果移动一态 \\
流水 trace & 写出 LOAD-use 停顿与随后分支误预测的 valid 变化 & 停顿保持，气泡插入，误预测清除 \\
\bottomrule
\end{longtable}

\section*{参考解答与要点}

\begin{longtable}{L{0.18\linewidth}L{0.46\linewidth}L{0.26\linewidth}}
\toprule
任务 & 参考解答 & 必须保留的检查点 \\
\midrule
有效位 & NOP 是有效指令且可能计入退休、异常和调试语义；气泡的 valid 为 0，所有架构写使能必须无效。 & 不用某个 opcode 代替所有权状态 \\
优先级 & 保存老指令异常 PC/原因，撤销年轻请求或标记其响应不可提交；流水转入异常入口，同时保证已接受事务不会失去归属。 & 老状态精确、年轻副作用不可见 \\
可变时延 & 至少保存 PC、目标寄存器、地址、宽度、符号扩展、请求 ID、异常和取消状态。 & 返回数据必须匹配原指令 \\
预测结构 & BHT 预测方向，BTB 给普通控制转移目标，RAS 给 return 目标；解析结果负责校验并恢复。 & 预测不等于提交 \\
计数器 trace & 弱跳转依次到强跳转、强跳转、弱跳转、弱不跳转、强不跳转。 & 强态需两次相反结果才翻转预测 \\
流水 trace & LOAD 继续进 MEM，ADD 前级保持并向 EX 插气泡；数据可旁路后恢复。分支误预测沿清除更年轻 valid 并重定向 PC。 & 老指令不被误清除 \\
\bottomrule
\end{longtable}
